% =============================================================================
%  capitulos/cap5_implementacao.tex
%  CAPÍTULO 5 — IMPLEMENTAÇÃO E DESENVOLVIMENTO
%
%  Objetivo: descrever o que foi de fato construído.
%  Pode haver diferenças em relação ao planejado — explique-as quando ocorrer.
%
%  NÃO cole todo o código-fonte aqui. Selecione trechos que:
%    - Resolvem um problema técnico não óbvio
%    - Implementam a lógica central do sistema
%    - Usam uma abordagem inovadora ou pouco comum
% =============================================================================

\chapter{Implementação e Desenvolvimento}
\label{cap:implementacao}

TODO: Parágrafo introdutório mencionando as principais tecnologias usadas,
como o capítulo está organizado e se houve alguma mudança relevante em
relação ao que foi planejado nos capítulos anteriores.


% -----------------------------------------------------------------------------
%  5.1 ARQUITETURA DO SISTEMA
%  Mostre como os componentes se comunicam entre si.
%  Ex.: React (frontend) → REST API (FastAPI) → PostgreSQL (banco)
%  Um diagrama de arquitetura aqui vale mais que mil palavras.
% -----------------------------------------------------------------------------
\section{Arquitetura do Sistema}
\label{sec:arquitetura}

TODO: Descreva as camadas e componentes do sistema e como eles se
comunicam. Mencione protocolos (HTTP/REST, WebSocket, gRPC...) e padrões
usados na comunicação entre as partes.

\begin{figure}[H]
  \caption{Arquitetura Geral do Sistema}
  \label{fig:arquitetura}
  \centering
  % \includegraphics[width=0.9\textwidth]{imagens/arquitetura.png}
  \fbox{\rule{0pt}{5cm}\rule{0.85\textwidth}{0pt}} % ← Placeholder, remova depois
  \fonte{Elaborado pelo autor.}
\end{figure}

TODO: Após o diagrama, descreva cada componente: o que faz, em que
tecnologia foi implementado e como se integra aos demais.


% -----------------------------------------------------------------------------
%  5.2 MODELAGEM DE DADOS
%  Apresente o DER (Diagrama Entidade-Relacionamento) se houver BD relacional.
%  Para NoSQL: descreva a estrutura dos documentos/coleções principais.
%  Remova esta seção se o projeto não usar banco de dados persistente.
% -----------------------------------------------------------------------------
\section{Modelagem de Dados}
\label{sec:modelagem-dados}

TODO: Apresente as principais entidades do sistema e seus relacionamentos.
A \cref{fig:der} exibe o Diagrama Entidade-Relacionamento do banco de dados.

\begin{figure}[H]
  \caption{Diagrama Entidade-Relacionamento (DER)}
  \label{fig:der}
  \centering
  % \includegraphics[width=0.9\textwidth]{imagens/der.png}
  \fbox{\rule{0pt}{7cm}\rule{0.9\textwidth}{0pt}} % ← Placeholder, remova depois
  \fonte{Elaborado pelo autor.}
\end{figure}

TODO: Após o DER, descreva as tabelas/entidades mais importantes,
seus atributos e as cardinalidades dos relacionamentos.


% -----------------------------------------------------------------------------
%  5.3 DESTAQUES DE IMPLEMENTAÇÃO
%  Subseções para funcionalidades relevantes ou tecnicamente interessantes.
%  Cada subseção deve ter: contexto → problema técnico → solução → código.
%  Use blocos de código CURTOS e FOCADOS — sem colar arquivos inteiros.
% -----------------------------------------------------------------------------
\section{Destaques de Implementação}
\label{sec:destaques}


% --- Destaque 1 ---
\subsection{TODO: Módulo ou Funcionalidade em Destaque 1}
\label{subsec:destaque1}

TODO: Explique o contexto: o que este módulo faz dentro do sistema?
Em seguida, descreva o problema técnico enfrentado e a solução adotada.
O \cref{cod:destaque1} apresenta o trecho responsável por [descreva o quê].

% Bloco de código — troque "Python" pela linguagem do seu projeto:
% Linguagens suportadas pelo listings: Python, Java, C, C++,
%   JavaScript, HTML, SQL, bash, PHP, Ruby, Swift, Kotlin, R...
\begin{lstlisting}[
  language=Python,
  caption={TODO: Descrição objetiva do que este código faz},
  label={cod:destaque1}
]
# TODO: Cole aqui um trecho curto e relevante do seu código.
# Foco: lógica central, não boilerplate.

def exemplo_funcao(parametro):
    """
    TODO: Docstring explicando a função.
    """
    # Comentários no código ajudam o leitor a acompanhar o raciocínio
    resultado = parametro * 2
    return resultado
\end{lstlisting}

TODO: Após o código, explique as linhas ou blocos mais importantes.
O leitor não deve precisar adivinhar o que acontece.


% --- Destaque 2 ---
\subsection{TODO: Módulo ou Funcionalidade em Destaque 2}
\label{subsec:destaque2}

TODO: Idem para outro trecho relevante do projeto.

% Exemplo de bloco JavaScript:
% \begin{lstlisting}[
%   language=JavaScript,
%   caption={Hook React para consumo da API de autenticação},
%   label={cod:auth-hook}
% ]
% import { useState } from 'react';
%
% const useAuth = () => {
%   const [user, setUser] = useState(null);
%
%   const login = async (credentials) => {
%     const response = await fetch('/api/auth/login', {
%       method: 'POST',
%       body: JSON.stringify(credentials),
%     });
%     const data = await response.json();
%     setUser(data.user);
%   };
%
%   return { user, login };
% };
% \end{lstlisting}


% --- Destaque 3 (descomente se necessário) ---
% \subsection{TODO: Módulo ou Funcionalidade em Destaque 3}
% \label{subsec:destaque3}
% TODO: Idem.


% -----------------------------------------------------------------------------
%  5.4 INTEGRAÇÃO ENTRE COMPONENTES (opcional)
%  Se a comunicação entre frontend, backend e banco de dados tiver algum
%  aspecto técnico relevante (autenticação JWT, WebSocket, cache...) vale
%  dedicar uma subseção para isso.
% -----------------------------------------------------------------------------
% \section{Integração entre Componentes}
% \label{sec:integracao}
% TODO: Descreva como as partes do sistema se comunicam em tempo de execução.
